\RequirePackage{plautopatch}
\documentclass[platex,dvipdfmx,a4paper]{jsarticle}

\usepackage{amsmath}
\usepackage{amssymb}
\newcommand{\hb}[2]{\frac{\sectionial #1}{\sectionial #2}}
\newcommand{\vectxy}[2]{\begin{pmatrix}#1\\#2\end{pmatrix}}
\newcommand{\vectxyz}[3]{\begin{pmatrix}#1\\#2\\#3\end{pmatrix}}
\newcommand{\vect}[1]{\boldsymbol{#1}}

\newcommand{\ssum}[3]{\sum_{#1=#2}^#3}

\title{物理数学I演習レポート2}
\author{嶋田 悠音}
\date{2026年5月10日}

\begin{document}
  \maketitle
  \subsection*{8.1}

  \subsubsection*{(i)}
  複素関数 $f(z) = \frac{1}{(z-a)(az-1)}$ を、原点を中心とする単位円 $\mathcal{C}$ （$z = e^{i\theta}, 0 \le \theta \le 2\pi$）上で積分することにより、定積分 $I = \int_{0}^{2\pi}\frac{d\theta}{1-2a \cos \theta+a^{2}}$ を求める。ただし $|a| < 1$ とする。

  単位円上では $z = e^{i\theta}$ より $dz = i e^{i\theta} d\theta = iz d\theta$、すなわち $d\theta = \frac{dz}{iz}$ である。また、$\cos \theta = \frac{1}{2}(z + z^{-1})$ である。
  これらを積分式に代入すると、分母は次のように変形できる：
  \begin{align*}
  1 - 2a \cos \theta + a^2 &= 1 - a(z + z^{-1}) + a^2 \\
  &= \frac{z - az^2 - a + a^2z}{z} \\
  &= \frac{-(z-a)(az-1)}{z}
  \end{align*}
  したがって、求める積分 $I$ は複素積分を用いて次のように書ける：
  \begin{align*}
  I &= \int_{|z|=1} \frac{z}{-(z-a)(az-1)} \frac{dz}{iz} = i \int_{|z|=1} \frac{1}{(z-a)(az-1)} dz \\
  &= i \int_{|z|=1} f(z) dz
  \end{align*}
  関数 $f(z) = \frac{1}{(z-a)(az-1)}$ の極は $z = a$ と $z = 1/a$ である。条件 $|a| < 1$ より、単位円の内部にある極は $z = a$ のみである。
  $z = a$ における留数 $\text{Res}(f, a)$ は、
  \begin{align*}
  \text{Res}(f, a) = \lim_{z \to a} (z-a) f(z) = \lim_{z \to a} \frac{1}{az-1} = \frac{1}{a^2-1}
  \end{align*}
  留数定理より、
  \begin{align*}
  I = i \cdot 2\pi i \cdot \text{Res}(f, a) = -2\pi \cdot \frac{1}{a^2-1} = \frac{2\pi}{1-a^2}
  \end{align*}

  \subsubsection*{(ii)}
  一般的な積分 $\int_{0}^{2\pi} f(\cos \theta, \sin \theta) d\theta$ を単位円 $\mathcal{C}$ に沿った複素積分で表す。
  $z = e^{i\theta}$ とおくと、$d\theta = \frac{dz}{iz}$ であり、以下の関係式を得る：
  \begin{align*}
  \cos \theta = \frac{z + z^{-1}}{2}, \quad \sin \theta = \frac{z - z^{-1}}{2i}
  \end{align*}
  これらを代入すると、積分は次のように表される：
  \begin{align*}
  \int_{0}^{2\pi} f(\cos \theta, \sin \theta) d\theta = \int_{\mathcal{C}} f\left( \frac{z + z^{-1}}{2}, \frac{z - z^{-1}}{2i} \right) \frac{dz}{iz}
  \end{align*}
  ここで、被積分関数を $g(z) = f\left( \frac{z + z^{-1}}{2}, \frac{z - z^{-1}}{2i} \right) \frac{1}{iz}$ と定義すれば、$\int_{\mathcal{C}} g(z) dz$ の形となる。

  \subsubsection*{(iii)}
  積分 $I_n = \int_{0}^{2\pi} \cos^n \theta \cos n\theta d\theta$ ($n \in \mathbb{N}$) を、原点を中心とする単位円 $\mathcal{C}$ 上の複素積分として計算する。

  単位円上の点 $z = e^{i\theta}$ において、$\cos \theta = \frac{z + z^{-1}}{2}$ および $d\theta = \frac{dz}{iz}$ である。また、$\cos n\theta = \text{Re}(z^n)$ であるが、計算を簡略化するため $\cos n\theta = \frac{z^n + z^{-n}}{2}$ として代入する。
  \begin{align*}
  I_n &= \int_{\mathcal{C}} \left( \frac{z + z^{-1}}{2} \right)^n \left( \frac{z^n + z^{-n}}{2} \right) \frac{dz}{iz} \\
  &= \frac{1}{i 2^{n+1}} \int_{\mathcal{C}} \frac{(z + z^{-1})^n (z^n + z^{-n})}{z} dz \\
  &= \frac{1}{i 2^{n+1}} \int_{\mathcal{C}} \frac{(z^2 + 1)^n (z^n + z^{-n})}{z^{n+1}} dz \\
  &= \frac{1}{i 2^{n+1}} \int_{\mathcal{C}} \left( \frac{(z^2 + 1)^n z^n}{z^{n+1}} + \frac{(z^2 + 1)^n z^{-n}}{z^{n+1}} \right) dz \\
  &= \frac{1}{i 2^{n+1}} \int_{\mathcal{C}} \left( \frac{(z^2 + 1)^n}{z} + \frac{(z^2 + 1)^n}{z^{2n+1}} \right) dz
  \end{align*}
  ここで、留数定理を用いる。被積分関数の各項について、単位円内部にある唯一の極 $z=0$ における留数を求める。

  1. 第一項 $\frac{(z^2 + 1)^n}{z}$ の $z=0$ における留数：
  $(z^2 + 1)^n$ の展開式における定数項（$z^0$ の係数）に等しい。
  二項定理より $(z^2 + 1)^n = \sum_{k=0}^n \binom{n}{k} (z^2)^k$ であるから、定数項は $k=0$ のとき、すなわち$1$ である。

  2. 第二項 $\frac{(z^2 + 1)^n}{z^{2n+1}}$ の $z=0$ における留数：
  $\frac{(z^2 + 1)^n}{z^2n}$の展開式における定数項、すなわち$(z^2 + 1)^n$ の展開式における $z^{2n}$ の係数に等しい。
  二項定理において $k=n$ のときであり、その係数は $1$ である。

  したがって、留数の総和は $1 + 1 = 2$ となる。留数定理より、
  \begin{align*}
  I_n &= \frac{1}{i 2^{n+1}} \cdot 2\pi i \cdot (1 + 1) \\
  &= \frac{2\pi \cdot 2}{2^{n+1}} \\
  &= \frac{4\pi}{2^{n+1}} \\
  &= \frac{\pi}{2^{n-1}}
  \end{align*}
  となる。
\end{document}